% Editorial revision of Sections 3 and 4.
% Retains the submitted models and measured results; technical issues requiring
% author decisions are tracked in system_paper_todos.md.
%==========================================================================
\section{Data Model and Problem Formulation}
\label{sec:model}

The motivating workload suggests that the representation needed to answer a
question can often be predicted before retrieving the image. We formalize the
assets, costs, and correctness estimates used to make this choice, then define
three optimization problems supported by \system.

\subsection{Multi-representation, multi-tier assets}

An image is an \emph{asset} $x$ with a set of available representations
$\R(x)\subseteq\R$. Our deployment uses five representations:
\[
\R=\{\textit{caption},\textit{thumb}_{32},\textit{jpeg}_{q1},
\textit{jpeg}_{q10},\textit{original}\}.
\]
A materialized instance $(x,r)$ has size $b_r(x)$ and resides at
$\ell(x,r)$ in the tier set
\[
\T=\{\textit{device},\textit{edge},\textit{cloud},\textit{archive}\}.
\]
Representations are derived from the original at ingestion. A
\emph{representation catalog} records each instance's location, size, and
storage pointer, and a pointer to caption text where applicable.
Planning reads catalog metadata, not images.

% TODO MODEL-01: P3 creates tier replicas but the above ell(x,r) permits
% only one tier; decide whether P1 selects (r,t) from an available-copy set.
\subsection{Tiers and cost vectors}

Retrieving representation $r$ of asset $x$ from tier $t$ has cost vector
\[
\mathbf{c}(x,r)=
\big(b_r(x),\;L_t(b_r(x)),\;C_t(b_r(x)),\;I_M(r)\big).
\]
Its components are transferred bytes, access latency, monetary cost, and VLM
inference time. We model access latency as $L_t(b)=\alpha_t+b/\beta_t$,
combining first-byte latency and transfer time. Monetary cost is
$C_t(b)=\kappa_t+\pi_t b$, combining per-request and per-byte charges.
Inference time $I_M(r)$ also depends on the representation: a caption is text,
whereas an image produces visual tokens. A weight vector $\mathbf{w}$ combines
these costs into a scalar objective. Applications may additionally impose hard
constraints, such as a retrieval deadline or spending limit.
% Evaluation scope (inference weight 0, archive not exercised) removed from the model by author, 10 Nov; inference scope stated in Sec. VI-B.

% TODO MODEL-02: Decide whether deadline D should include inference and
% optimizer overhead, rather than only L_t(b).
\subsection{Queries and correctness}

A query $Q=(q,x,M,\mathbf{w},D)$ asks question $q$ about asset $x$ using VLM
$M$, under weights $\mathbf{w}$ and deadline $D$. Let
$Y_r(q,x)\in\{0,1\}$ indicate whether $M$ answers correctly using
representation $r$. This outcome is not known when the optimizer makes its
decision. Instead, it estimates
\begin{equation}
 p_r(q)=\Pr\big[Y_r(q,x)=1\mid\phi(q)\big],
 \label{eq:success}
\end{equation}
where $\phi(q)$ contains features extracted from the question and, optionally,
precomputed catalog statistics. The estimator does not read a representation
of $x$. The probability describes expected success over the workload, not a
guarantee for an individual image.

\subsection{Three optimization problems}

\begin{definition}[\textbf{P1: Per-query selection}]
Given $Q$ and estimates $\widehat p_r$, choose
\begin{equation}
 r^*(Q)=\arg\max_{r\in\R(x):\,L_{\ell(x,r)}(b_r(x))\le D}
 \big[\widehat p_r(q)-\lambda\,\mathbf{w}^\top\mathbf{c}(x,r)\big].
 \label{eq:p1}
\end{equation}
\end{definition}

The objective values a correct answer at one unit and trades expected
correctness against weighted access and inference costs. Selecting $r$ incurs
the cost of that representation alone, rather than the cumulative cost of all
smaller representations. This distinguishes direct selection from a cascade.

\begin{definition}[\textbf{P2: Workload allocation}]
Given a batch $\mathcal{Q}$ and a budget $B$ on scalar cost
$c=\mathbf{w}^\top\mathbf{c}$, choose one representation $r(Q)$ for each query
to maximize $\sum_{Q\in\mathcal{Q}}\widehat p_{r(Q)}$, subject to
$\sum_{Q\in\mathcal{Q}}c(x_Q,r(Q))\le B$.
\end{definition}

\begin{definition}[\textbf{P3: Representation placement}]
Given a workload distribution $\mathcal{W}$, tier capacities $K_t$, and
per-tier cost models, choose which asset--representation instances to
materialize on each tier. The placement maximizes expected P1 utility for
$\mathcal{W}$ subject to the storage capacity of each tier.
\end{definition}

\textbf{P1} handles queries as they arrive. \textbf{P2} allocates a shared budget across a known
batch, as when a response mission has a fixed backhaul allowance. \textbf{P3} is an
offline physical-design problem, analogous to materialized-view and index
selection~\cite{harinarayan1996cubes,chaudhuri1997autoadmin}. We do not solve P3
optimally; \S\ref{sec:design} gives a greedy placement advisor, which we evaluate in a
two-tier simulation. The same
correctness estimates used at query time also score candidate placements.

% TODO MODEL-01: Reinstate a precise set/capacity formulation for P3 once
% replica semantics and the asset dimension are resolved.
\subsection{Guarantees}

The quality of P1 depends on estimation error. The following bound makes this
dependence explicit.

\begin{proposition}[Selection regret]
\label{prop:regret}
For query $Q$, let
$\epsilon(Q)=\max_r|\widehat p_r(q)-p_r(q)|$ over feasible representations.
Let $r^*$ maximize the true objective in~\eqref{eq:p1}, and let $\widehat r$
maximize the estimated objective over the same feasible set. Then the utility
lost by choosing $\widehat r$ is at most $2\epsilon(Q)$. The expected regret
is therefore at most $2\,\mathbb{E}_Q[\epsilon(Q)]$.
\end{proposition}

\begin{proof}
Add and subtract the estimated utilities of $r^*$ and $\widehat r$. The
difference between their estimated utilities is non-positive because
$\widehat r$ maximizes that objective. Each of the remaining two terms is
bounded by $\epsilon(Q)$; the cost terms cancel because the same costs are
used in both objectives. Taking expectations gives the second claim.
\end{proof}

The guarantee relates selection regret to the estimator's error on the
operating workload. It does not bound error uniformly across unseen query
distributions or establish a finite-sample guarantee for calibration. We
therefore use disjoint populations for estimator training, calibration, and
operating-point selection (\S\ref{sec:design}).

P2 is a multiple-choice knapsack problem (MCKP). Each query is a class and
each representation is an item, with value $\widehat p_r$ and cost $c$.
Although MCKP is NP-hard, its LP relaxation can be solved by considering
incremental value per cost on the upper convex hull of each class. An optimal
basic LP solution has at most one fractional class~\cite{sinha1979mckp}.
Provided the cheapest allocation is feasible, rounding that class down loses
at most one expected correct answer relative to the integer optimum, since
$0\le\widehat p_r\le1$. This gives a useful bound for large batches.

% TODO MODEL-03: Verify the LP/rounding statement and assumptions against
% the cited result and distinguish it from the bucketed-DP implementation.
%==========================================================================
\section{\system Design}
\label{sec:design}

\begin{figure*}[t]
\centering
\begin{tikzpicture}[
 font=\footnotesize,
 box/.style={draw,rounded corners,align=center,minimum height=9mm,text width=25mm,fill=white},
 store/.style={draw,align=center,minimum height=9mm,text width=25mm,fill=gray!10},
 arrow/.style={-{Latex[length=1.6mm]},thick},
 darrow/.style={-{Latex[length=1.6mm]},dashed},
 node distance=7mm and 6mm]
\node[box] (q) {Query $(q,x,M)$\\weights $\mathbf{w}$, $D$, $B$};
\node[box,right=of q] (feat) {Feature extraction\\$\phi(q)$ (text only)};
\node[box,right=of feat] (est) {Success estimators $g_{M,r}$\\and calibrators $\psi_{M,r}$};
\node[box,right=of est] (opt) {Optimizer\\P1 per query, P2 batch};
\node[box,right=of opt] (exec) {Executor: read $(x,r^*)$\\and invoke $M$ once};
\node[box,below=of q] (ing) {Ingestion and\\materializer};
\node[box,below=of feat] (prof) {Offline profiler\\and trainer};
\node[box,below=of est] (place) {Placement\\advisor (P3)};
\node[box,below=of opt] (cat) {Representation\\catalog};
\node[store,below=of exec] (tiers) {Tiered storage: device,\\edge, cloud, archive};
\draw[arrow] (q)--(feat);
\draw[arrow] (feat)--(est);
\draw[arrow] (est)--(opt);
\draw[arrow] (opt)--(exec);
\draw[arrow] (cat)--(opt);
\draw[arrow] (exec)--(tiers);
\draw[darrow] (ing)--(prof);
\draw[darrow] (prof)--(est);
\draw[darrow] (est)--(place);
\draw[darrow] (place)--(cat);
\draw[darrow] (cat)--(tiers);
\draw[darrow] (ing.south) -- ++(0,-4mm) -| (tiers.south);
\end{tikzpicture}
\caption{\system architecture. The solid path performs query-time selection
using the question and catalog metadata; only the executor reads a
representation. Dashed arrows show offline ingestion, profiling, and
placement.}
\label{fig:arch}
\end{figure*}

Figure~\ref{fig:arch} shows the online and offline components of \system.
At query time, the system extracts text features, estimates the probability
of a correct answer for each available representation, and combines those
estimates with cataloged access costs. The executor then retrieves the chosen
representation and invokes the target VLM once. Separately, ingestion creates
representations, profiling trains and calibrates the estimators, and a
placement advisor determines which tiers should hold them.

\subsection{Query interface and catalog}

\system exposes representation selection as an access-path choice for a VLM
operator. Consider a semantic query:
\begin{center}
\footnotesize
\begin{tabular}{l}
\texttt{SELECT a.id, ASK(a.image, 'How many buildings}\\
\texttt{\ \ are flooded?', 'qwen3-vl') FROM aerial a}\\
\texttt{WHERE a.sector = 4}\\
\texttt{WITH (accuracy\_floor = 0.95, deadline = 1s)}
\end{tabular}
\end{center}

Here \texttt{a.image} denotes an asset, rather than its bytes. The query
planner leaves the representation passed to \texttt{ASK} unresolved. For each
qualifying asset, \system uses the question and catalog metadata to choose a
representation before retrieving the image. The executor can group these
reads by tier, while P2 can allocate a shared budget across the qualifying
queries. This decision can also be inserted into an existing semantic query
engine~\cite{patel2024lotus,liu2025palimpzest}: \system chooses the input
representation, leaving operator ordering and model selection to the host
engine.

The representation catalog has schema
\begin{center}
\footnotesize
\begin{tabular}{l}
\texttt{Rep(asset\_id, rep, tier, uri, bytes,}\\
\texttt{\ \ derived\_from, created\_at, caption\_text\_uri)}
\end{tabular}
\end{center}
and stores one row per materialized instance. Tier parameters
$(\alpha_t,\beta_t,\kappa_t,\pi_t)$ are maintained separately. Changes in
prices or measured bandwidth update cost estimates without retraining the
success estimators. Each target VLM also has a profile containing its
representation estimators, calibrators, validated operating point, and label
provenance. Since planning requires only metadata and model profiles, the
catalog can be replicated at the edge or on devices.

% TODO DESIGN-01: The example's accuracy_floor=0.95 must be aligned with
% the actual relative validation floor of 0.95*A_fix.
% TODO MODEL-01: If replicas are permitted, Rep's key must include tier.
\subsection{Ingestion and materialization}

The materializer generates representations where the original is first
available. It creates thumbnails and JPEG previews by re-encoding the image,
and uses a captioning model to generate a caption. Captioning is performed
once per asset, so its creation cost is amortized over subsequent queries.
Our online cost accounting excludes that one-time cost, which we did not
measure; we also evaluate \system without captions (\S\ref{sec:robust}). The materializer records the
size and location of each resulting representation in the catalog.

\subsection{Offline profiling}

For each target VLM $M$, profiling runs the model on labeled questions at
every representation and records the correctness outcome $Y_r$. This is the
only stage in the design that deliberately evaluates multiple representations
for the same question. The resulting labels are used to train the success
estimators. Each outer fold's training questions are split by group into 70\%
for the estimators, 15\% for calibration, and 15\% for operating-point
selection; on FloodNet this leaves 1,156 labeled questions per fold. Fitting
all five estimators and calibrators takes 0.5\,s on FloodNet and 41\,s on
12,000 VQA-v2 questions on one laptop CPU (Apple M4).
% src: VOILA_code scripts/run_all_experiments.py L1457-1474 (70/15/15 split, leakage asserts);
% generated_icde/planner_overhead.txt (scripts/planner_overhead.py, own re-implementation, same hyperparameters)

\subsection{Success estimators}

The input to each estimator is a feature vector $\phi(q)$ built from the
question text. It contains sparse TF--IDF features and structural indicators,
including counting expressions, yes/no auxiliaries, numerals, negation,
superlatives, and question length. For each model--representation pair
$(M,r)$, a gradient-boosted estimator $g_{M,r}$ predicts a raw success score.
The estimators learn from the model's observed correctness labels; they do not
inspect image content at query time.

The estimator family is replaceable. With calibration and the optimizer held
fixed, gradient boosting, logistic regression, ridge regression, and an MLP
differ by at most 1.6 accuracy points for \mPix on VQA-v2, and a fine-tuned
BERT regressor in the same pipeline does not consistently improve the
accuracy--cost trade-off (both in \S\ref{sec:robust}). We use gradient
boosting because it has low training and prediction overhead in this setting.
The interface also permits features computed at ingestion and recorded in the
catalog, such as representation sizes. We do not rely on such features for the
reported question-text-only results.
% TODO E3: Evaluate catalog-statistic features; do not imply a completed result.

\subsection{Per-representation calibration}

Raw estimator scores are not directly comparable: a score of 0.6 for a
caption need not imply the same success probability as 0.6 for the original.
Because the optimizer trades estimated correctness against cost, these scores
must be calibrated. For each $(M,r)$, \system fits an isotonic map
$\psi_{M,r}$ on a population disjoint from estimator training (the 15\%
calibration split), producing
\[
\widehat p_r(q)=\psi_{M,r}\big(g_{M,r}(\phi(q))\big).
\]
We calibrate each representation separately and do not require correctness to
increase with representation size; Table~\ref{tab:fixed} shows why such a
constraint would be inappropriate.

\subsection{Per-query optimizer}

Algorithm~\ref{alg:p1} implements P1. Using catalog metadata, it excludes
representations that are unavailable for the asset or violate the retrieval
deadline. It evaluates the lightweight estimators for the remaining
candidates, computes their expected utilities, and selects the maximizer. The
online path reads one representation and invokes the VLM once, rather than
retrieving progressively larger inputs. Planning a query takes 0.64\,ms
(median; 7\,$\mu$s per query when batched) on one laptop CPU core, which is
small next to the access latencies of \S\ref{sec:objectives}.
% src: generated_icde/planner_overhead.txt (scripts/planner_overhead.py)

The operating-point parameter $\lambda$ is selected on validation data,
separate from both training and calibration. Let
$A^{\mathrm{val}}_{\mathrm{fix}}$ be the accuracy of the best fixed
representation on that population. Among candidate penalties achieving at
least $0.95A^{\mathrm{val}}_{\mathrm{fix}}$, \system selects the one that
transfers the fewest bytes. This is a relative validation accuracy target,
not a per-query guarantee; applications may choose a stricter trade-off.

Unlike sequential escalation, P1 compares the full utility of each available
representation. An earlier incremental-gain heuristic remains in the
experimental ablations, but does not implement the P1 objective.

\begin{algorithm}[t]
\caption{Per-query representation selection (P1)}
\label{alg:p1}
\begin{algorithmic}[1]
\Require query $(q,x,M)$, weights $\mathbf{w}$, deadline $D$, penalty $\lambda$
\State $F\gets\{r\in\R(x): L_{\ell(x,r)}(b_r(x))\le D\}$ \Comment{catalog only}
\If{$F=\emptyset$}
  \State $F\gets\{\arg\min_{r\in\R(x)}L_{\ell(x,r)}(b_r(x))\}$
  \Comment{best effort; deadline unmet}
\EndIf
\ForAll{$r\in F$}
  \State $\widehat p_r\gets\psi_{M,r}(g_{M,r}(\phi(q)))$
  \State $u_r\gets\widehat p_r-\lambda\,\mathbf{w}^\top\mathbf{c}(x,r)$
\EndFor
\State $r^*\gets\arg\max_{r\in F}u_r$
\State Read $(x,r^*)$ from $\ell(x,r^*)$; \Return $M(q,(x,r^*))$
\end{algorithmic}
\end{algorithm}

% TODO DESIGN-02: Decide if best-effort deadline violation is permitted
% by the API; otherwise explicitly return infeasible.
\subsection{Workload optimizer}

For a known batch with a shared budget, \system solves P2 as an MCKP: each
query contributes one class, with a candidate item for each representation.
We implement two allocators. \emph{Bucketed dynamic programming} rounds costs
up to 1\,KiB buckets and retains Pareto-optimal (cost, value) states. Rounding
up makes the returned allocation feasible under the unrounded budget, though
it need not be optimal for the original continuous-cost instance. In our
measurements, a 100-query batch is solved in under a second.
\emph{Marginal-value allocation} is a greedy alternative: it initially assigns
each query its cheapest representation, then applies upgrades in decreasing
order of predicted accuracy gain per additional unit cost, provided they fit
within the remaining budget.
% src: VOILA_code scripts/percom_experiments.py mckp_dynamic_program, greedy_upgrades

Both methods plan over a batch; single queries use P1.

\subsection{Placement advisor}
\label{sec:advisor}

P3 uses the same correctness estimates to guide representation placement.
For a candidate placement, the placement advisor replays a sample of the workload
through the P1 decision rule, restricted to representations available at the
modeled tiers. It measures the expected utility under that placement. Starting
from representations stored in the cloud, it greedily copies the
asset--representation instance with the largest popularity-weighted utility
gain per added storage byte to a faster tier, stopping at the tier's capacity.
Popularity weights come from the expected workload. The cloud copy is
retained, and P1 reads the fastest available copy. A copy can improve utility
by lowering a tier-dependent cost or by making a representation feasible
under a deadline. With the byte objective we evaluate, cost does not depend on
the tier, so all of the placement advisor's measured gains come from deadline
feasibility. Placement evaluation uses the success estimators rather than
invoking VLMs, so the placement advisor scores every (image, representation) candidate,
7,225 on FloodNet, offline (\S\ref{sec:placement-eval}).

% TODO MODEL-01: The P1 equations and catalog must support selecting among
% multiple copies of the same representation for this advisor to be exact.
% TODO DESIGN-03: Reconcile cloud-at-ingestion initial placement with
% on-device/edge materialization and originals aging into archive.
\subsection{Maintenance}

\textbf{New representations.} Since estimators are defined per
model--representation pair, adding a representation requires profiling and
training its estimator and calibrator; existing representation estimators do
not need to be retrained.

\textbf{New models.} Estimators trained on an earlier model can provide a
starting point for a newer one, with recalibration as new labels become
available. Within the Qwen family, zero-shot transfer is 3.4--4.5 accuracy
points below a target-trained estimator on VQA-v2 and matches its accuracy on
GQA (\S\ref{sec:robust}).

\textbf{Unseen question forms.} Question-text estimators are less reliable on
templates absent during profiling (\S\ref{sec:robust}). For workloads drawn
from a fixed set of question templates, as in FloodNet, the model profile
records the normalized question strings (lowercase, punctuation removed) seen
during profiling. A query whose string is not recorded reads the best fixed
representation on the profiling data until labels for its template are
collected. Because matching is exact, a rephrased profiled question is also
treated as new. That forgoes savings on the query but avoids the held-out
template failure. A feature-distance test does not work: every held-out
FloodNet template has TF--IDF cosine similarity of at least 0.8 to a profiled
one (\S\ref{sec:robust}). Exact matching does not suit open-ended workloads.
On VQA-v2, GQA, and TextVQA, 69--82\% of test questions have a string never
seen in training, so the fallback would disable selection for most queries.
We do not use it on those workloads.
% src: generated_icde/template_coverage.txt (scripts/template_coverage.py); generated_icde/ood/
% DESIGN-08 (was DESIGN-04) resolved 10 Nov; false negatives cannot occur with exact matching,
% false positives (rephrasings) are unmeasured because FloodNet has none.
